% ==========================================
% BAB VI EVALUASI  (versi terisi hasil validasi 7 kasus)
% Berkas chapter tanpa preamble; disisipkan ke naskah skripsi.
% Membutuhkan definisi kolom L dan C{..} serta paket booktabs, tabularx,
% array, enumitem (sudah tersedia pada preamble skripsi Anda).
% ==========================================
\cleardoublepage
\chapter{EVALUASI}
\label{chap:evaluasi}
Bab ini memaparkan evaluasi terhadap sistem yang dibangun melalui dua sudut
pandang yang saling melengkapi. Pertama, evaluasi kuantitatif terhadap model
klasifikasi \emph{Machine Learning} yang menjadi inti prediksi risiko hipertensi,
yang diukur menggunakan metrik-metrik terstandar. Kedua, evaluasi kualitatif
terhadap dua komponen sistem yang bersifat interpretatif, yaitu komponen
\emph{Explainable Artificial Intelligence} (SHAP) dan komponen \emph{Large Language
Model} (LLM), yang keluarannya tidak dapat sepenuhnya diukur oleh metrik otomatis
sehingga memerlukan penilaian pakar. Evaluasi kuantitatif merupakan rekapitulasi
atas hasil pemodelan yang telah dilaksanakan pada Bab~V, sedangkan evaluasi
kualitatif dilakukan melalui pendekatan \emph{expert validation} oleh seorang
dokter umum sebagai \emph{domain expert}.

\section{Evaluasi Kuantitatif (ML)}
\label{sec:evaluasi-kuantitatif}
Bagian ini merekapitulasi hasil evaluasi kuantitatif model \emph{Machine Learning}
yang telah dilaksanakan secara rinci pada tahap pemodelan di Bab~V. Tujuan
rekapitulasi ini adalah menegaskan model terbaik yang menjadi mesin prediksi
sistem beserta kinerjanya, sebelum pembahasan berpindah ke evaluasi kualitatif
komponen interpretatif. Uraian lengkap mengenai pembagian data, penyetelan
\emph{hyperparameter}, optimasi ambang keputusan, hingga perbandingan antarmodel
dapat ditelusuri kembali pada Bab~V.

Pada tahap pemodelan, empat algoritma dibandingkan, yaitu \emph{Logistic
Regression}, \emph{Random Forest}, \emph{XGBoost}, dan \emph{CatBoost}. Keempatnya
disetel menggunakan \texttt{GridSearchCV} dengan fungsi objektif ROC-AUC dan
validasi silang tiga lipatan, kemudian titik operasinya ditentukan melalui
\emph{Youden Index}. Pemilihan model akhir tidak hanya bertumpu pada satu metrik,
melainkan pada gabungan pertimbangan kemampuan membedakan kelas (ROC-AUC),
kemampuan menangkap pasien berisiko (\emph{recall}), kemudahan penafsiran,
dukungan analisis kontribusi fitur berbasis SHAP, serta kemudahan penerapan pada
sistem. Berdasarkan pertimbangan tersebut, \textbf{XGBoost} dipilih sebagai model
akhir karena memiliki kemampuan membedakan kelas yang setara dengan model terbaik,
\emph{recall} tertinggi pada ambang bawaan, dukungan penjelasan SHAP yang matang,
dan penerapan yang ringan.

Kinerja model terbaik XGBoost pada data uji dirangkum pada
Tabel~\ref{tab:rekap-ml}, disajikan pada dua titik operasi, yaitu ambang bawaan
0,5 dan ambang \emph{Youden} yang menjadi titik operasi final sistem.

\begin{table}[htbp]
\centering
\caption{Rekapitulasi kinerja model terbaik (XGBoost) pada data uji}
\label{tab:rekap-ml}
\begin{tabular}{@{}lcc@{}}
\toprule
\textbf{Metrik} & \textbf{Ambang 0,5} & \textbf{Ambang Youden} \\
\midrule
ROC-AUC       & \multicolumn{2}{c}{0,7795} \\
PR-AUC        & \multicolumn{2}{c}{0,5795} \\
Recall        & 0,7186 & 0,7062 \\
Specificity   & 0,7036 & 0,7192 \\
F1-\emph{score} & 0,5901 & 0,5921 \\
Akurasi       & 0,7080 & 0,7154 \\
\bottomrule
\end{tabular}
\end{table}

Nilai ROC-AUC dan PR-AUC dihitung dari probabilitas sehingga tidak berubah oleh
pemilihan ambang. Dengan ROC-AUC sebesar 0,7795 dan \emph{recall} sebesar 0,7186
pada ambang 0,5 (0,7062 pada ambang Youden), model XGBoost memenuhi kriteria
keberhasilan yang ditetapkan di awal penelitian, yaitu ROC-AUC minimal 0,75 dan
\emph{recall} minimal 0,70. Model ini kemudian disimpan bersama nilai ambang
\emph{Youden} dan urutan nama fitur, lalu digunakan sebagai mesin prediksi yang
keluarannya menjadi masukan bagi komponen SHAP dan LLM. Kualitas kedua komponen
interpretatif itulah yang dievaluasi secara kualitatif pada bagian berikut.

\section{Evaluasi Kualitatif (XAI + LLM)}
\label{sec:evaluasi-kualitatif}
Berbeda dengan model klasifikasi yang dapat diukur secara kuantitatif, komponen
SHAP dan LLM menghasilkan keluaran interpretatif yang kualitasnya tidak dapat
sepenuhnya ditangkap oleh metrik otomatis sehingga memerlukan penilaian pakar.
Evaluasi kualitatif dilakukan melalui pendekatan \emph{expert validation} oleh
seorang dokter umum sebagai \emph{domain expert}.

\subsection{Metode Evaluasi}
\label{sec:metode-evaluasi}

\subsubsection{Strategi Evaluasi}
Strategi evaluasi yang digunakan adalah validasi ahli (\emph{expert validation})
oleh satu orang dokter umum yang berperan sebagai \emph{domain expert}. Pendekatan
ini dipilih, dan bukan pengukuran kesepakatan antar-penilai (\emph{inter-rater
agreement}), karena penelitian melibatkan satu penilai tunggal. Metrik kesepakatan
antar-penilai seperti Cohen's Kappa atau Aiken's V mensyaratkan minimal dua penilai
dan tidak dapat dihitung secara sahih untuk penilai tunggal; oleh sebab itu fokus
evaluasi diarahkan pada validasi kualitas keluaran sistem oleh ahli, dengan analisis
deskriptif dan tematik terhadap hasil penilaian.

Objek evaluasi terdiri atas dua hal yang dinilai secara terpisah menggunakan dua
instrumen yang berbeda: (1) penjelasan kontribusi fitur yang dihasilkan SHAP terhadap
prediksi risiko hipertensi, dan (2) narasi personal edukatif yang dihasilkan LLM
berdasarkan hasil SHAP. Pemisahan instrumen dilakukan karena kedua komponen memiliki
objek penilaian yang berbeda: SHAP dinilai sebagai penjelasan interpretatif model,
sedangkan LLM dinilai sebagai teks komunikatif untuk pasien.

\subsubsection{Dasar Pemilihan Instrumen}
\label{sec:dasar-instrumen}
Seluruh aspek penilaian pada kedua instrumen diturunkan dari literatur bereputasi
agar tidak ada aspek yang ditentukan secara subjektif. Aspek untuk komponen SHAP
diadaptasi dari Salih dkk. (2024) yang terbit di \emph{Artificial Intelligence Review}
(jurnal Q1), sebuah tinjauan sistematis atas pendekatan evaluasi XAI pada aplikasi
klinis kardiovaskular. Tinjauan tersebut mengkategorikan evaluasi XAI menjadi lima
jalur (\emph{human-grounded}, \emph{application-grounded}/pakar domain,
\emph{proxy-grounded}, \emph{literature-grounded}, dan \emph{guideline-grounded}),
serta mendefinisikan kriteria kualitas penjelasan klinis seperti relevansi klinis,
\emph{plausibility}, \emph{complexity}/keterpahaman, dan kepercayaan sebagai tujuan
akhir evaluasi pakar.

Aspek untuk komponen LLM diadaptasi dari Tam dkk. (2024) yang terbit di
\emph{npj Digital Medicine} (jurnal Q1), yang mengusulkan kerangka QUEST untuk
evaluasi LLM di bidang kesehatan berdasarkan tinjauan atas 142 studi. QUEST terdiri
atas lima prinsip beserta dimensinya: \emph{Quality of Information} (antara lain
\emph{accuracy}, \emph{agreement}, \emph{consistency}, \emph{relevance},
\emph{usefulness}), \emph{Understanding and Reasoning}, \emph{Expression Style and
Persona} (\emph{clarity}, \emph{empathy}), \emph{Safety and Harm} (\emph{bias},
\emph{harm}, \emph{fabrication/falsification}), dan \emph{Trust and Confidence}.
Kerangka ini juga menegaskan penggunaan skala Likert 1--5 dengan definisi tiap skor
sebagai format penilaian yang lazim, sehingga menjadi dasar desain skala pada kedua
instrumen penelitian ini.

Tabel~\ref{tab:penelusuran-shap} dan Tabel~\ref{tab:penelusuran-llm} menunjukkan
penelusuran setiap aspek instrumen ke dimensi asalnya pada literatur.

\begin{table}[htbp]
\centering
\caption{Penelusuran aspek instrumen SHAP ke Salih dkk. (2024)}
\label{tab:penelusuran-shap}
\begin{tabularx}{\textwidth}{@{}p{0.28\textwidth} L@{}}
\toprule
\textbf{Aspek instrumen} & \textbf{Dimensi/kriteria pada Salih dkk. (2024)} \\
\midrule
Relevansi fitur        & \emph{Clinical relevance} --- penjelasan selaras dengan opini dokter dan mendukung penalaran klinis; dikuantifikasi melalui persentase tumpang tindih fitur XAI dengan pilihan dokter. \\
Arah kontribusi fitur  & Pembahasan kolinearitas --- prediktor yang saling berkorelasi (mis. tekanan darah, merokok, diabetes) dapat membuat XAI menghasilkan penjelasan tidak realistis/bias, sehingga arah kontribusi perlu diverifikasi ahli. \\
Kesesuaian medis       & \emph{Plausibility} dan jalur \emph{literature-grounded} --- penjelasan sejalan dengan ekspektasi pakar dan pengetahuan medis mapan. \\
Kemudahan dipahami     & \emph{Complexity}/\emph{understandability} --- kompleksitas dan waktu yang dibutuhkan untuk memahami penjelasan. \\
Kepercayaan            & \emph{Trust} --- keterlibatan pakar bertujuan menilai apakah penjelasan layak dipercaya, digeneralisasi, dan diterapkan. \\
\bottomrule
\end{tabularx}
\end{table}

\begin{table}[htbp]
\centering
\caption{Penelusuran aspek instrumen LLM ke kerangka QUEST, Tam dkk. (2024)}
\label{tab:penelusuran-llm}
\begin{tabularx}{\textwidth}{@{}p{0.30\textwidth} L@{}}
\toprule
\textbf{Aspek instrumen} & \textbf{Prinsip/dimensi pada QUEST (Tam dkk., 2024)} \\
\midrule
Faktualitas medis          & \emph{Quality of Information}: \emph{accuracy} + \emph{agreement}. \\
Konsistensi terhadap SHAP  & \emph{Quality of Information}: \emph{consistency} (kesesuaian isi narasi dengan bukti SHAP sebagai rujukan). \\
Koherensi \& kejelasan     & \emph{Expression Style and Persona}: \emph{clarity}; koherensi naratif sebagai aspek kunci evaluasi. \\
Kualitas edukasi           & \emph{Quality of Information}: \emph{usefulness} + \emph{relevance}; edukasi pasien merupakan kategori aplikasi utama. \\
Bebas halusinasi           & \emph{Safety and Harm}: \emph{fabrication/falsification}. \\
Keamanan (potensi bahaya)  & \emph{Safety and Harm}: \emph{harm}. \\
\bottomrule
\end{tabularx}
\end{table}

\subsubsection{Objek dan Prosedur Validasi}
\label{sec:prosedur}
Validasi dilakukan oleh dr. Roos Meri Silaen (Puskesmas Mulyamekar, Dinas Kesehatan
Kabupaten Purwakarta) pada tanggal 8 Juli 2026 atas tujuh kasus representatif yang
dipilih dari data uji. Pemilihan kasus mempertimbangkan keterwakilan kelas dan hasil
prediksi, yaitu mencakup kasus berisiko rendah hingga berisiko tinggi serta kombinasi
faktor risiko yang beragam, agar dokter menilai perilaku sistem pada beragam kondisi,
bukan hanya pada kasus yang mudah. Profil ketujuh kasus disajikan pada
Tabel~\ref{tab:profil-kasus}. Untuk setiap kasus, dokter disajikan: (a) ringkasan data
pasien; (b) hasil prediksi model beserta probabilitasnya; (c) visualisasi kontribusi
fitur SHAP; dan (d) narasi personal yang dihasilkan LLM.

\begin{table}[htbp]
\centering
\caption{Profil tujuh kasus yang divalidasi}
\label{tab:profil-kasus}
\small
\begin{tabularx}{\textwidth}{@{}c p{0.18\textwidth} c c L@{}}
\toprule
\textbf{Kasus} & \textbf{Profil} & \textbf{Jenis kelamin} & \textbf{Usia} & \textbf{Faktor risiko menonjol} \\
\midrule
1 & Atlet amatir           & Laki-laki  & 21 & Tanpa faktor risiko; kualitas tidur sangat baik \\
2 & Mahasiswa tingkat akhir & Perempuan & 22 & Kualitas tidur buruk; gangguan tidur sering \\
3 & Pekerja korporat       & Laki-laki  & 34 & Merokok; kolesterol tinggi \\
4 & Ibu rumah tangga       & Perempuan  & 48 & Diabetes; gangguan tidur sering \\
5 & Supir truk             & Laki-laki  & 45 & Merokok; kolesterol tinggi; kualitas tidur sangat buruk \\
6 & Pensiunan aktif        & Perempuan  & 65 & Tanpa komorbid; usia lanjut \\
7 & Komorbid lengkap       & Laki-laki  & 52 & Merokok; diabetes; kolesterol tinggi; IMT tinggi \\
\bottomrule
\end{tabularx}
\end{table}

Prosedur validasi ditempuh dalam langkah berikut:
\begin{enumerate}[leftmargin=1.4em,itemsep=2pt,topsep=2pt]
    \item Dokter menerima penjelasan mengenai tujuan penelitian, komponen sistem, dan
    petunjuk pengisian kedua lembar validasi.
    \item Untuk setiap kasus, dokter meninjau keluaran SHAP dan mengisi Lembar Validasi
    SHAP pada kelima aspeknya.
    \item Untuk kasus yang sama, dokter meninjau narasi LLM dan mengisi Lembar Validasi
    LLM pada keenam aspeknya.
    \item Dokter mengisi kolom komentar untuk memberikan justifikasi kualitatif atas
    skor yang diberikan dan mencatat temuan spesifik.
    \item Skor pada seluruh kasus direkap untuk dianalisis.
\end{enumerate}

Kedua instrumen menggunakan skala Likert 1--5 dengan orientasi seragam ``semakin
tinggi semakin baik''. Untuk aspek yang secara alami bermuatan negatif (halusinasi
dan potensi bahaya), konstruk yang dinilai dirumuskan secara positif (``bebas
halusinasi'' dan ``keamanan/bebas potensi bahaya'') sehingga skor 5 selalu menandakan
kualitas terbaik. Perumusan ini menjaga agar rerata antaraspek dapat diinterpretasikan
secara konsisten tanpa \emph{reverse coding}.

\subsubsection{Metode Analisis Hasil Evaluasi}
\label{sec:analisis}
Pada penelitian ini, analisis dilakukan secara deskriptif dan tematik. Untuk analisis deskriptif, setiap aspek dihitung rerata skor Likert
lintas kasus beserta ukuran sebarannya (median dan modus). Rerata tiap aspek kemudian
diinterpretasikan menggunakan skala interval lima kategori sebagaimana pada
Tabel~\ref{tab:interpretasi}. Rerata keseluruhan per komponen (SHAP dan LLM) juga
dihitung sebagai indikator ringkas kualitas komponen.

\begin{table}[htbp]
\centering
\caption{Skala interpretasi rerata skor Likert}
\label{tab:interpretasi}
\begin{tabular}{@{}C{3.2cm}C{3.6cm}C{3.6cm}@{}}
\toprule
\textbf{Rentang rerata} & \textbf{Kategori} & \textbf{Interpretasi validitas} \\
\midrule
1{,}00 -- 1{,}80 & Sangat kurang & Tidak valid \\
1{,}81 -- 2{,}60 & Kurang        & Kurang valid \\
2{,}61 -- 3{,}40 & Cukup         & Cukup valid \\
3{,}41 -- 4{,}20 & Baik          & Valid \\
4{,}21 -- 5{,}00 & Sangat baik   & Sangat valid \\
\bottomrule
\end{tabular}
\end{table}

Tahap selanjutnya adalah pada analisis tematik. Komentar dokter dianalisis secara kualitatif untuk
mengidentifikasi pola temuan, misalnya jenis ketidaksesuaian fitur, arah kontribusi
yang meragukan, atau bentuk halusinasi pada narasi. Analisis tematik ini melengkapi
angka Likert dengan konteks klinis dan menjadi dasar rekomendasi perbaikan sistem.

\subsection{Hasil Evaluasi}
\label{sec:hasil-evaluasi}
Validasi dilaksanakan atas tujuh kasus dan menghasilkan 35 penilaian untuk komponen
SHAP (tujuh kasus dikali lima aspek) serta 42 penilaian untuk komponen LLM (tujuh kasus
dikali enam aspek). Skor rinci per kasus disajikan pada Tabel~\ref{tab:perkasus-shap}
dan Tabel~\ref{tab:perkasus-llm}, sedangkan ringkasan rerata, median, modus, dan
kategori setiap aspek disajikan pada Tabel~\ref{tab:hasil-shap} dan
Tabel~\ref{tab:hasil-llm}.

\begin{table}[htbp]
\centering
\caption{Skor validasi komponen SHAP per kasus}
\label{tab:perkasus-shap}
\small
\begin{tabular}{@{}clcccccccc@{}}
\toprule
\textbf{No} & \textbf{Aspek} & \textbf{K1} & \textbf{K2} & \textbf{K3} & \textbf{K4} & \textbf{K5} & \textbf{K6} & \textbf{K7} & \textbf{Rerata} \\
\midrule
1 & Relevansi fitur        & 5 & 5 & 4 & 5 & 4 & 4 & 3 & 4,29 \\
2 & Arah kontribusi fitur  & 2 & 3 & 4 & 3 & 3 & 2 & 3 & 2,86 \\
3 & Kesesuaian medis       & 3 & 3 & 4 & 3 & 4 & 2 & 3 & 3,14 \\
4 & Kemudahan dipahami     & 4 & 4 & 5 & 4 & 4 & 5 & 5 & 4,43 \\
5 & Kepercayaan            & 4 & 4 & 4 & 3 & 3 & 2 & 3 & 3,29 \\
\bottomrule
\end{tabular}
\end{table}

\begin{table}[htbp]
\centering
\caption{Skor validasi komponen LLM per kasus}
\label{tab:perkasus-llm}
\small
\begin{tabular}{@{}clcccccccc@{}}
\toprule
\textbf{No} & \textbf{Aspek} & \textbf{K1} & \textbf{K2} & \textbf{K3} & \textbf{K4} & \textbf{K5} & \textbf{K6} & \textbf{K7} & \textbf{Rerata} \\
\midrule
1 & Faktualitas medis          & 3 & 5 & 5 & 4 & 4 & 2 & 3 & 3,71 \\
2 & Konsistensi terhadap SHAP  & 3 & 4 & 4 & 4 & 3 & 3 & 4 & 3,57 \\
3 & Koherensi \& kejelasan     & 5 & 5 & 5 & 4 & 3 & 3 & 4 & 4,14 \\
4 & Kualitas edukasi           & 3 & 4 & 4 & 5 & 4 & 2 & 3 & 3,57 \\
5 & Bebas halusinasi           & 3 & 4 & 5 & 4 & 5 & 3 & 4 & 4,00 \\
6 & Keamanan (potensi bahaya)  & 5 & 5 & 5 & 5 & 5 & 4 & 3 & 4,57 \\
\bottomrule
\end{tabular}
\end{table}

\begin{table}[htbp]
\centering
\caption{Hasil validasi komponen SHAP}
\label{tab:hasil-shap}
\begin{tabular}{@{}clcccc@{}}
\toprule
\textbf{No} & \textbf{Aspek} & \textbf{Rerata} & \textbf{Median} & \textbf{Modus} & \textbf{Kategori} \\
\midrule
1 & Relevansi fitur       & 4,29 & 4 & 4; 5 & Sangat baik \\
2 & Arah kontribusi fitur & 2,86 & 3 & 3    & Cukup \\
3 & Kesesuaian medis      & 3,14 & 3 & 3    & Cukup \\
4 & Kemudahan dipahami    & 4,43 & 4 & 4    & Sangat baik \\
5 & Kepercayaan           & 3,29 & 3 & 3; 4 & Cukup \\
\midrule
\multicolumn{2}{@{}l}{\textbf{Rerata keseluruhan}} & 3,60 & & & Baik \\
\bottomrule
\end{tabular}
\end{table}

\begin{table}[htbp]
\centering
\caption{Hasil validasi komponen LLM}
\label{tab:hasil-llm}
\begin{tabular}{@{}clcccc@{}}
\toprule
\textbf{No} & \textbf{Aspek} & \textbf{Rerata} & \textbf{Median} & \textbf{Modus} & \textbf{Kategori} \\
\midrule
1 & Faktualitas medis          & 3,71 & 4 & 3; 4; 5 & Baik \\
2 & Konsistensi terhadap SHAP  & 3,57 & 4 & 4       & Baik \\
3 & Koherensi \& kejelasan     & 4,14 & 4 & 5       & Baik \\
4 & Kualitas edukasi           & 3,57 & 4 & 4       & Baik \\
5 & Bebas halusinasi           & 4,00 & 4 & 4       & Baik \\
6 & Keamanan (potensi bahaya)  & 4,57 & 5 & 5       & Sangat baik \\
\midrule
\multicolumn{2}{@{}l}{\textbf{Rerata keseluruhan}} & 3,93 & & & Baik \\
\bottomrule
\end{tabular}
\end{table}

Selain skor, komentar dokter memberikan konteks kualitatif atas penilaian. Pada
komponen SHAP, temuan yang paling menonjol adalah ketidaksesuaian arah kontribusi pada
beberapa fitur berdampak kecil, khususnya merokok, status tidak
merokok, dan jenis kelamin, yang pada seluruh kasus dinilai masih dapat diterima karena
magnitudonya berada di bawah 3\%. Pada komponen LLM, temuan yang menonjol mencakup
saran yang tidak sesuai dengan kondisi pasien (Kasus 1), kesalahan logika kausal
terkait indeks massa tubuh (Kasus 4), ketidakkonsistenan saran berhenti merokok
terhadap arah kontribusi SHAP (Kasus 5), serta penilaian bobot diabetes yang dianggap
terlalu kecil (Kasus 7). Catatan lintas kasus yang menonjol adalah keterkaitan antara
galat arah pada SHAP dan ketidakkonsistenan narasi LLM, yang diuraikan lebih lanjut
pada Subbab~\ref{sec:pembahasan-evaluasi}.

\subsection{Pembahasan Hasil Evaluasi}
\label{sec:pembahasan-evaluasi}
Pembahasan menghubungkan hasil penilaian dokter dengan kualitas komponen XAI dan
narasi LLM, serta implikasinya bagi kelayakan sistem.

\subsubsection{Interpretasi Kualitas Komponen SHAP}
Hasil validasi SHAP pada Tabel~\ref{tab:hasil-shap} menunjukkan rerata keseluruhan
3,60 yang tergolong valid. Dua aspek memperoleh kategori sangat baik, yaitu relevansi
fitur (4,29) dan kemudahan dipahami (4,43). Nilai tersebut menunjukkan bahwa fitur yang
ditandai SHAP sebagai paling berkontribusi umumnya memang relevan secara klinis menurut
penilaian dokter, dan bahwa penyajian kontribusi fitur mudah dipahami tanpa keahlian
teknis tambahan. Temuan ini selaras dengan kriteria \emph{clinical relevance} dan
\emph{understandability} pada Salih dkk. (2024).

Tiga aspek lain memperoleh kategori cukup, yaitu arah kontribusi fitur (2,86),
kesesuaian medis (3,14), dan kepercayaan (3,29). Aspek arah kontribusi fitur menjadi
titik terlemah komponen SHAP. Komentar dokter secara konsisten menandai adanya fitur
yang arah kontribusinya tidak sesuai dengan hubungan yang dikenal secara medis. Pada
beberapa kasus, fitur gangguan tidur dan merokok yang seharusnya menaikkan risiko
justru ditampilkan sebagai penurun risiko, dan pada Kasus 6 status tidak merokok yang
seharusnya menurunkan risiko malah ditampilkan menambah risiko sebesar 1,3\%.
Ketidaksesuaian arah tersebut konsisten dengan pembahasan Salih dkk. (2024) mengenai
pengaruh kolinearitas antarprediktor kardiovaskular yang dapat membuat metode XAI
menghasilkan penjelasan yang tidak realistis.

Meskipun demikian, dokter menilai ketidaksesuaian itu masih dapat diterima karena
magnitudo kontribusinya sangat kecil, yaitu umumnya di bawah 3\%, sehingga tidak
menggeser kesimpulan tingkat risiko. Penilaian bersyarat ini penting untuk dicatat:
validitas arah kontribusi SHAP tergolong baik pada fitur berdampak besar, tetapi
mengandung galat arah yang sistematis pada fitur berdampak kecil. Galat itulah yang
menyeret turun skor kesesuaian medis dan kepercayaan, dan tampak paling jelas pada
Kasus 6 yang memuat beberapa galat arah sekaligus serta memperoleh skor kepercayaan
terendah, yaitu 2.

\subsubsection{Interpretasi Kualitas Narasi LLM}
Komponen LLM memperoleh rerata keseluruhan 3,93 yang tergolong valid, lebih tinggi
daripada komponen SHAP. Aspek keamanan memperoleh nilai tertinggi (4,57; sangat baik),
diikuti koherensi dan kejelasan (4,14). Hal ini menunjukkan bahwa narasi yang
dihasilkan aman untuk disampaikan kepada pasien serta tersusun runtut dan mudah
dipahami. Aspek faktualitas medis (3,71) dan bebas halusinasi (4,00) juga tergolong
valid, yang menandakan narasi umumnya benar secara medis dan jarang memuat informasi
yang dibuat-buat.

Dua aspek dengan nilai terendah adalah konsistensi terhadap SHAP (3,57) dan kualitas
edukasi (3,57). Komentar dokter memperlihatkan penyebab yang spesifik dan berulang.
Pada Kasus 1, narasi menyarankan pasien memperbaiki kualitas tidur padahal data pasien
menunjukkan kualitas tidur sudah sangat baik. Pada Kasus 4 terdapat kesalahan logika
kausal, yaitu narasi menyatakan bahwa perbaikan kualitas tidur ``dapat meningkatkan
indeks massa tubuh'', padahal maksud yang benar adalah menurunkannya. Pada Kasus 6,
keterkaitan antara interpretasi kontribusi indeks massa tubuh dan saran yang diberikan
dinilai kurang padu sehingga aspek faktualitas hanya memperoleh skor 2. Pada Kasus 7,
kontribusi diabetes dinilai terlalu kecil, padahal secara klinis resistensi insulin
pada diabetes meningkatkan aktivitas saraf simpatis dan menyempitkan pembuluh darah
sehingga tergolong faktor moderat. Sebagian ketidaksesuaian ini berkaitan dengan
strategi \emph{prompt engineering} yang belum sepenuhnya memaksa narasi mempertahankan
makna dan proporsi kontribusi dari hasil SHAP.

\subsubsection{Keterkaitan Antarkomponen dan Implikasi Sistem}
Hasil kedua komponen saling berkaitan pada alur XAI$\rightarrow$LLM. Kasus 5
memperlihatkan keterkaitan ini dengan jelas: hasil SHAP menandai merokok sebagai
penurun risiko, yaitu arah yang keliru secara medis, sementara narasi LLM menyarankan
pasien berhenti merokok. Saran tersebut benar secara medis, tetapi tidak konsisten
dengan arah kontribusi yang ditampilkan SHAP, sehingga menimbulkan ketegangan antara
kesetiaan terhadap sumber (\emph{faithfulness}) dan kebenaran medis. Pola ini
menunjukkan bahwa sebagian persoalan konsistensi pada komponen LLM berakar pada galat
arah di komponen SHAP. Perbaikan pada komponen SHAP di hulu, dengan demikian,
berpotensi langsung meningkatkan konsistensi narasi di hilir.

Secara keseluruhan, kedua komponen memperoleh kategori valid, yaitu SHAP 3,60 dan
LLM 3,93. Temuan ini menunjukkan bahwa rangkaian interpretasi SHAP dan narasi LLM layak
digunakan sebagai pendukung komunikasi risiko hipertensi kepada pasien, namun belum
memadai untuk berdiri sendiri sebagai alat pengambilan keputusan klinis. Galat arah
pada fitur berdampak kecil dan sejumlah ketidakkonsistenan narasi perlu diperbaiki
sebelum sistem diterapkan pada konteks yang lebih luas.

\subsubsection{Keterbatasan}
Evaluasi ini memiliki sejumlah keterbatasan. Pertama, penilaian dilakukan oleh satu
orang dokter umum sehingga hasilnya tidak dapat digeneralisasi secara luas dan tidak
memungkinkan pengukuran kesepakatan antar-penilai. Kedua, jumlah kasus yang dievaluasi
terbatas pada tujuh profil. Ketiga, penilaian ahli bersifat subjektif; pada Kasus 1
dokter mencatat bahwa profil tersebut merupakan profil pertama yang dinilai sehingga
terdapat kemungkinan efek urutan pada awal penilaian. Keempat, sebagian penilaian aspek
arah kontribusi bersifat bersyarat, yaitu ketidaksesuaian arah dinilai dapat diterima
sepanjang magnitudonya kecil, sehingga validitas aspek tersebut perlu dimaknai dengan
mempertimbangkan syarat ini.

\subsubsection{Rekomendasi Perbaikan Sistem}
Berdasarkan temuan di atas, terdapat beberapa rekomendasi perbaikan. Pertama, galat
arah pada fitur berdampak kecil dapat ditangani dengan penyaringan atau penggabungan
kontribusi di bawah ambang tertentu agar tidak ditampilkan sebagai arah yang
menyesatkan. Kedua, strategi \emph{prompt engineering} pada komponen LLM perlu
diperkuat agar narasi mempertahankan arah dan proporsi kontribusi dari hasil SHAP serta
menghindari kesalahan logika kausal seperti pada Kasus 4. Ketiga, penelitian lanjutan
disarankan menambah jumlah penilai untuk memungkinkan pengukuran kesepakatan
antar-penilai serta memperluas cakupan kasus yang dievaluasi.

% ==========================================
% Catatan panduan penulisan (dari template) tetap dipertahankan sebagai komentar.
% ==========================================
% \section{Evaluasi Kesalahan Penulisan yang Sering Terjadi}
% \subsection{Penggunaan Kata "di mana" atau "dimana"}
% Banyak yang menuliskan kata "di mana" atau "dimana" sebagai pengganti kata "which" dalam bahasa Inggris.
% Padahal, penggunaan kata "di mana" atau "dimana" tidak tepat dalam konteks tersebut. Demikian juga untuk kata serupa, misalnya "yang mana".
% Kata "di mana" atau "dimana" ini harus diganti dengan kata lain, seperti "dengan", "tempat", "yang", dan sebagainya tergantung kalimatnya.
% Penjelasan lengkap dapat dilihat pada \autocite{BPBI}.